\subsection{Продвинутая работа с камерой}

Мы уже использовали камеру для поиска маркеров. Теперь разберемся, как работать с видеопотоком более гибко.

\subsubsection*{Класс VideoStream}
Для получения <<сырого>> видеопотока (без автоматического декодирования в OpenCV формат, что иногда нужно для оптимизации) используется класс \texttt{VideoStream}.

\begin{codefile}[python]{python/examples/camera_examples/camera_stream.py}
\end{codefile}

Этот код запускает передачу видео по Wi-Fi. Вы можете открыть этот поток в плеере VLC (адрес обычно \texttt{rtsp://192.168.4.1:8554/front}). Как и вся эта часть курса, RTSP-трансляция требует реального дрона: браузерный симулятор не отдает поток наружу, открыть его в VLC не получится.

\subsubsection*{Захват кадров и фотографирование}
Чтобы сделать фото, нужно получить текущий кадр и сохранить его в файл с помощью \texttt{cv2.imwrite}.

\begin{taskbox}[Фотограф]
Напишите программу, которая показывает видео с дрона на экране. При нажатии клавиши \texttt{S} (Save) текущий кадр должен сохраняться в файл \texttt{photo\_N.jpg}, где N — порядковый номер фото.
\end{taskbox}

\subsubsection*{Калибровка камеры}
Для точных измерений (например, расстояния до маркера) нужно знать оптические параметры камеры: матрицу камеры и коэффициенты дисторсии.

Процесс калибровки:
1. Распечатать шахматную доску (Chessboard).
2. Сделать серию снимков доски под разными углами.
3. Запустить скрипт калибровки (\texttt{camera\_calibration.py}), который вычислит параметры и сохранит их в файл (обычно \texttt{.yml} или \texttt{.xml}).

\begin{infobox}[Зачем это нужно?]
Без калибровки прямые линии на изображении могут казаться изогнутыми (эффект <<рыбьего глаза>>), что приведет к ошибкам в навигации.
\end{infobox}
